\section{Introduction}
\label{sec:intro}

The paradigm of large-scale pre-training followed by task-specific fine-tuning has revolutionized deep learning, producing a vast library of specialized neural networks sharing a common initialization. To utilize these specialized models without the linear scaling in memory and compute associated with classic ensembling or Mixture-of-Experts (MoE) \cite{shazeer2017outrageously, fedus2022switch}, the community has turned to weight-space model merging. Methods such as model soups \cite{wortsman2022model} and task arithmetic \cite{ilharco2022editing} interpolate the parameters of fine-tuned networks directly, executing multi-task inference in a single forward pass with zero parameter or memory overhead.

Despite their efficiency, standard static model-merging techniques suffer from a fundamental limitation: they apply a single, fixed interpolation weight across all input samples. This static compromise is highly susceptible to interference under severe task conflict, where task-specific parameter directions clash. To address this, recent work has proposed \emph{dynamic} model merging \cite{gu2024zero, yadav2024resolving}, which computes sample-specific merging coefficients on the fly using input-dependent routers. 

A central question in dynamic model merging is the spatial granularity of the routing policy. Should we route at a global, model-wide level, or should each layer in the deep hierarchy possess its own dynamic routing policy? Intuition from representation learning suggests that because different layers in deeply hierarchical networks extract highly distinct semantic abstractions (e.g., low-level features in early layers versus task-specific, high-level concepts in deep layers), different layers should benefit from specialized routing weights. 

However, a highly influential recent theoretical result \cite{anonymous} presented a mathematical proof of ``Layer-Averaging Collapse'' (or rank-1 collapse) to assert that layer-wise routing is entirely redundant. The authors claimed that in any dynamic model-merging system, learned layer-wise routing coefficient trajectories must become perfectly collinear across all layers, effectively collapsing the routing space to a single, global dimension. This theoretical claim has strongly steered the community: recent architectures have abandoned layer-specific routing policies in favor of computationally simpler, global, single-layer routers, citing the mathematical redundancy of layer-wise parameterization.

In this work, we adopt a critical methodological perspective to audit this claim. We argue that the theoretical proof of ``Layer-Averaging Collapse'' is contingent on highly idealized assumptions---specifically, perfectly collinear base representations and strictly contractive linear Jacobians. In realistic multi-task settings where hierarchical deep networks must learn heterogeneous, cross-domain tasks, these assumptions are strongly violated.

To isolate and probe these routing dynamics with absolute precision and transparency, we establish a mathematically rigorous, physical empirical framework directly on Split-MNIST digit subsets. We train task-specific expert weights on subsets of MNIST digits: Task 0: digits 0 and 1; Task 1: digits 2 and 3; Task 2: digits 4 and 5; Task 3: digits 6 and 7. We then implement a dynamic weight-space model merging pipeline on two representative architectures: a \textbf{DeepMLP-12} with $L=12$ layers (Transformer block equivalent), and a \textbf{TinyCNN-4} with $L=4$ layers (residual stage equivalent). Under a tight few-shot calibration budget, we evaluate performance across three distinct multi-task suites representing varying degrees of semantic conflict.

Under this physical framework, we construct the \textbf{Batch-Averaged Layer-wise Coefficient Matrix} $A \in \mathbb{R}^{L \times K}$ (where $L$ is the number of layers and $K$ is the number of expert models) and probe its true dimensionality using Singular Value Decomposition (SVD). Additionally, we compute the inter-layer pairwise cosine similarity matrix to map the directional alignment of routing decisions along the network's depth.

Our rigorous, multi-seed physical evaluation deconstructs the Layer-Averaging Collapse claim and yields several key methodological insights:
\begin{itemize}
    \item \textbf{Deconstruction of Rank-1 Collapse:} Under our physical Cross-Domain task suite, the SVD Collinearity Ratio $\rho_{collinear}$ drops to \textbf{0.4987} for DeepMLP-12 and \textbf{0.5673} for TinyCNN-4. This demonstrates that in high-conflict settings, learned routing trajectories occupy a multi-dimensional subspace rather than collapsing to a rank-1 line, deconstructing prior theoretical claims.
    \item \textbf{Emergence of Depth-Specialized Policies:} Inter-layer cosine similarity maps reveal that while low-conflict suites yield highly uniform routing across layers, cross-domain conflict forces the network to specialize its routing into distinct block-diagonal structures, aligning with hierarchical feature extraction.
    \item \textbf{The Capacity-Variance Trade-off:} We expose a fundamental trade-off: under low-data budgets, the lower parameter capacity of global routers acts as an implicit regularizer, whereas our layer-wise router successfully scales to higher budgets to outperform simpler baselines as variance is controlled.
    \item \textbf{Critical Role of Regularization:} We highlight how router regularization governs generalization. Omitting weight decay stabilizes calibration under low task conflict, enabling unregularized layer-wise routers to find highly expressive local interpolations.
\end{itemize}

By exposing these structural limits of representational collinearity, we demonstrate that layer-wise dynamic model merging is not a redundant over-parameterization, but an expressive, depth-specialized framework capable of achieving superior multi-task consensus under high task conflict.
